% Appendix B: sampled decisions, further analysis.
% Proofs: Theorem fixednoise (statement in the main text), Propositions coupling,
% commonmass, boundedrange, kl, frontier, Theorems guards, residualrace, residualself.
% None of these constructions is implemented on the GPU; the synthetic experiment is a
% CPU work count (register entry synthetic).
\section{Sampled decisions: further analysis}\label{app:sampling}
The engine path of Section~\ref{sec:decisions} verifies sampled drafts with a fixed noise field, which turns every sampled decision into a race. This appendix proves the theorem behind that choice and develops what the choice avoids: the standard acceptance rule evaluated with enclosures, residual races and adaptive verification depth. Every statement takes enclosures of the logits as input and does not depend on the mechanism that produced them, so it applies to transport as well as to the int8 copy. None of these alternatives is implemented on the GPU. Sampling before the scores are known, an approximate mode, a demand frontier and a CPU experiment that shows why the alternatives were not built are in the research notes.

\subsection{Fixed-noise verification and the rest of the random-number contract}\label{app:fixednoise}
\begin{proof}[Proof of Theorem~\ref{thm:fixednoise}]
By induction on positions, every accepted token equals the argmax at a prefix equal to $Y_{<t}$, so it is $Y_t$; at the first mismatch the prefix is $Y_{<t}$ and the emitted token is $Y_t$. The drafter changes only how many positions one target pass verifies, so the output is the same for every drafter and depth policy. Given $Y_{<t}$, the field $G_{t,\cdot}$ is independent of $z_t(Y_{<t})$, so $Y_t\mid Y_{<t}\sim\softmax(z_t(Y_{<t})/T)$ by the Gumbel-max identity.
\end{proof}
The theorem treats $z_t$ as a function of the prefix alone, which the engine does not guarantee (Section~\ref{sec:same-shape} and Appendix~\ref{app:divergence}); the contract of Section~\ref{sec:decisions} is scoped accordingly.

A keyed field is what gives pathwise equality with the dense race: in the exact rational emulation, feeding the candidates noise from a stream consumed in evaluation order changed the winner in 247 of 300 cases~\ev{precision}.

Two further conditions complete the random-number contract of Section~\ref{sec:decisions}. Equality in distribution, as opposed to seeded equality, needs i.i.d.\ noise with the correct marginal that is independent of the proposal draw, the acceptance uniform and the scores; a residual race (Appendix~\ref{app:residual}) also needs independence from the acceptance event. And a hash with $2^{32}$ levels makes the law exact only for the idealized continuous field; the stock seeded sampler shares that approximation.

\paragraph{Truncated sampling}\label{app:truncated} The Gumbel-max identity has a top-$k$ form that samples $k$ tokens without replacement~\cite{kool2019}. Top-$k$ sampling applies the ladder to membership (drop $i$ when $hi_i$ is below the $k$-th largest $lo$) and races over the certified set. With noise keyed by token, as in the full-vocabulary race, this gives equality in distribution with top-$k$ sampling from the reference scores and, once the engine's tie semantics at the $k$-th position are matched, pathwise equality with a token-keyed dense race over the same set. It does not give seeded equality with SGLang's sampler at the pin.

That sampler's seeded top-$k$ and top-$p$ path sorts the probabilities, keys each token's Gumbel noise by the token's rank among the sorted probabilities rather than by its index, and maps the winner back to a token through the sort's permutation, so the noise a token receives depends on its rank~\cite{sglangsampler}. Seeded equality with that path therefore needs, beyond membership, the complete order within the certified set: enclosures of the FP32 probabilities the sampler sorts that are strictly separated for every pair of members whose ranks could otherwise differ, and \code{torch.sort}'s order where probabilities are equal. The pinned sampler also refuses min-$p$ with a seed. Its full-vocabulary seeded path keys the noise by token index, and that is the case the certified race of Section~\ref{sec:decisions} covers.

Other truncations differ in what they need. Min-$p$ admits $i$ when $z_i\ge z_{\max}+T\log(\text{min-}p)$, which intervals decide without a partition function. Top-$p$ admission depends on the mass above each token relative to the partition function $Z$, so it needs partition bounds and is usually cheaper to hand to the ordinary path. Fused samplers take the other route and make the dense path cheap: FlashInfer samples top-$k$, top-$p$ and min-$p$ without sorting and implements chain speculative sampling~\cite{flashinfer,flashinfersamplingblog}, and SonicSampler fuses the post-logit sampling and verification kernels over a bounded top-128 pool and describes top-$p$ under that pool as theoretically lossy~\cite{sonic}. The closest precedent for sampling from bounded scores, Mussmann, Levy and Ermon's lazy Gumbel sampling, draws noise only for a retrieved top set and, with an approximate top set, widens the cutoff by the approximation error~\cite{mussmann2017,mussmann2016}; the certified race differs in bounding every vocabulary entry with a dense cheap pass and in keying the noise by token, which gives pathwise equality with a dense race keyed the same way, not with the stock sampler's rank-keyed truncated path.

\subsection{Acceptance under fixed-noise coupling}\label{app:coupling}
Fixed-noise verification of a sampled draft accepts less often than the standard rule. The proposition says by how much.
\begin{proposition}[Acceptance under fixed-noise coupling]\label{prop:coupling}
Suppose that the drafter proposes the token $\argmax_i(\log q_i+G_i)$ and that the target accepts it exactly when it equals $\argmax_i(\log p_i+G_i)$. Then
\begin{gather*}
P(\mathrm{accept})=\sum_i\Bigl(\sum_j\max\bigl(\tfrac{p_j}{p_i},\tfrac{q_j}{q_i}\bigr)\Bigr)^{-1},\\
\tfrac12\sum_i\min(p_i,q_i)\le\sum_i\frac{p_iq_i}{p_i+q_i}\le P(\mathrm{accept})\le\sum_i\min(p_i,q_i).
\end{gather*}
\end{proposition}
\begin{proof}
Both choose $i$ if and only if $G_i>\max_{j\ne i}(G_j+m_j)$ with $e^{m_j}=\max(p_j/p_i,q_j/q_i)$, an event of probability $(1+\sum_{j\ne i}e^{m_j})^{-1}$. The bounds follow from $\max(x,y)\le x+y$, from $\max(x,y)\ge x$ and $\max(x,y)\ge y$ summed over $j$, and from $p+q\le2\max(p,q)$.
\end{proof}
Per position, fixed-noise acceptance of a sampled draft is therefore at most the standard rule's rate $1-\TV(p,q)$ and at least half of it, and the two agree when $p=q$. For SGLang's default point-mass drafts the question does not arise: the stock rule accepts a draft $d$ with probability $p(d)$, and so does the keyed race. Whether a cheaper certificate pays for any lost acceptance with sampled drafts is an empirical question, and the acceptance under fixed-noise coupling against standard verification on real drafts has not been measured~\pend{geometry-accept}.

\subsection{The standard acceptance rule with enclosures}\label{app:guards}\label{sec:guards}\label{sec:standard-rule}
At the pin, SGLang's sampled speculative verification uses point-mass drafts by default: a drafted token is accepted when a uniform coin falls below the target probability mass of the siblings tried so far, and a rejection samples the target restricted to the untried tokens by inverse CDF (code reading). In general, let $p$ be the target law at the reached prefix and $q$ the \emph{actual} proposal law, conditional on the information used to produce the proposal. A softmax buffer from a different selector is not $q$, and a deterministic proposal is a point mass. The standard one-step correction accepts $X\sim q$ with probability $\min(1,p_X/q_X)$ and otherwise samples $(p-q)_+$ normalized; the accepted mass at $i$ is $\min(p_i,q_i)$ and the residual contributes $(p_i-q_i)_+$, which sum to $p_i$~\cite{specdecode,chen2023specsampling}. SpecTr and block verification extend the rule to several drafts and to whole blocks~\cite{spectr,blockverify}. A certificate may change how this comparison is evaluated, never the law used in it.

For one reached proposal $X=x$ write $w_x=e^{z_x}$ and $Z=\sum_ie^{z_i}$, and draw $U\in[0,1)$ with the reference convention. For $q_x>0$, acceptance is equivalent to
\begin{equation}Uq_xZ<w_x,\label{eq:accept}\end{equation}
and the $\min(1,\cdot)$ need not be evaluated because $U<1$ already handles $p_x\ge q_x$.

\begin{theorem}[Acceptance guards with enclosures]\label{thm:guards}\label{thm:guard}
Let $w_x\in[N^-,N^+]$ and $Z\in[Z^-,Z^+]$ with $N^->0$, let $0<q_x\le1$ be the actual proposal mass and $U\in[0,1)$. Then $Uq_xZ^+<N^-$ implies acceptance, $Uq_xZ^-\ge N^+$ implies rejection, the two never both hold, and replacing an enclosure by a tighter one cannot reverse a certified decision. For $U$ uniform and independent of the enclosures, neither guard fires with probability
\[P_?=\min\bigl(1,N^+/(q_xZ^-)\bigr)-\min\bigl(1,N^-/(q_xZ^+)\bigr).\]
If $N^-=N^+=w_x$ and every entry outside a re-scored set $R$ has a mass enclosure whose upper and lower ends differ by a factor at most $\kappa$, then with $\pi=\sum_{i\notin R}p_i$,
\[P_?\le\frac{p_x}{q_x}\cdot\frac{(\kappa-1)\pi}{1-\pi(1-1/\kappa)}.\]
\end{theorem}
\begin{proof}
The implications are transitivity with~\eqref{eq:accept}; both firing would give $N^+\le Uq_xZ^-\le Uq_xZ^+<N^-$. A tighter enclosure stays on the same side of the true values. The unresolved uniforms form $[N^-/(q_xZ^+),N^+/(q_xZ^-))\cap[0,1)$. For the bound, $P_?\le(w_x/q_x)(Z^+-Z^-)/(Z^+Z^-)$; $Z^+-Z^-\le\sum_{i\notin R}e^{z_i}(\kappa-1)$ because each lower mass is at most $e^{z_i}$; and $Z^-\ge Z-\sum_{i\notin R}e^{z_i}(1-1/\kappa)$ because each lower mass is at least $e^{z_i}/\kappa$. Divide by $Z$.
\end{proof}
Transported tile masses (Theorem~\ref{thm:transport}) and exponentiated intervals from the int8 copy (Theorem~\ref{thm:envelope}) both supply $Z^\pm$. With the int8 copy $\kappa=e^{2\hat\beta}$ times the outward rounding of the exponential; at $\beta=1$ ($\kappa\approx7.4$), keeping $P_?$ below 10\% at $p_x/q_x\approx1$ requires the unre-scored mass $\pi$ to be below about 1.5\%. After an adaptive refinement chosen using $U$ the probability statement no longer applies, although the guards remain sound pointwise. The measured partition widths on real draft states are in Appendix~\ref{app:transport-sampled}; for transport they leave $P_?$ near one.

The guards decide acceptance, and that is where they stop being enough. A rejection still requires a sample from $(p-q)_+$, and knowing that rejection occurred does not supply one. Under a stock-token contract the stock inverse-CDF residual would need every logit of the position after each rejection, so a cheap acceptance decision would still be followed by a pass over the whole head. This is why the engine path uses fixed-noise verification, which needs neither a partition function nor a residual sampler, and why the guards and the residual races below are analysis, not built.

\subsection{Residual and bonus races}\label{app:residual}\label{sec:residual}
When all drafted tokens are accepted, a bonus token needs a target sample; after a rejection the residual needs one. Two facts about the engine and one construction reduce the problem. First, SGLang's default sampled verification uses point-mass proposals (Appendix~\ref{app:guards}); with $q$ a point mass at $x$, the residual is the target restricted to the untried tokens, a race over $V\setminus\{x\}$ that needs no partition function. Second, under fixed-noise verification there is no residual at all. Third, for a proposal with small support there is an exact bounded race.

Let $Q=\{i:q_i>0\}$. With unnormalized target weights $w_i$ and partition $Z$, the residual can be sampled from the weights
\begin{equation}
r_i=(w_i-Zq_i)_+,\qquad\frac{r_i}{\sum_jr_j}=\frac{(p_i-q_i)_+}{\sum_j(p_j-q_j)_+},\label{eq:residualweights}
\end{equation}
with $r_i=w_i$ outside $Q$. Given $Z^-\le Z\le Z^+$, positive-part monotonicity gives $(w_i-Z^+q_i)_+\le r_i\le(w_i-Z^-q_i)_+$ inside $Q$.

\begin{theorem}[Bounded residual race]\label{thm:residualrace}
Assume valid enclosures of the target weights and of $Z$, the actual proposal law, a nonzero residual, and a priority field $\pi_i=e^{G_i}>0$ that is fixed and independent of the proposal draw, the acceptance uniform and the rejection event. A refinement procedure that returns only an interval-certified ordered maximizer of $r_i\pi_i$, or completes the dense race, returns the same token as the dense race. With i.i.d.\ standard Gumbel $G_i$ it samples the residual law~\eqref{eq:residualweights} exactly.
\end{theorem}
\begin{proof}
The enclosures of $r_i$ multiplied by $\pi_i>0$ enclose the priorities. A certified maximizer dominates every possible competitor, including the actual one, and a completed evaluation yields the dense ordered maximizer. The distributional statement is Gumbel-max applied to the nonzero residual weights.
\end{proof}

\begin{theorem}[Residual races from weight enclosures]\label{thm:residualself}
Replace the enclosures above by $r_i\in[(N^-_i-Z^+q_i)_+,(N^+_i-Z^-q_i)_+]$ for $q_i>0$ and $r_i\in[N_i^-,N_i^+]$ for $q_i=0$. The conclusions of Theorem~\ref{thm:residualrace} hold. For a point-mass proposal at $x$, $r_x=0$ as soon as $Z^-\ge N^+_x$, and the residual is a race over $V\setminus\{x\}$ that needs no partition bound. SGLang's target-only rule (accept sibling $j$ if $UZ<\sum_{l\le j}w_{x_l}$; otherwise sample the target restricted to the untried tokens) is certified by the guards with a cumulative numerator and an exclusion race.
\end{theorem}
\begin{proof}
Positive-part monotonicity gives the enclosures; the rest is the proof of Theorem~\ref{thm:residualrace}. For the point mass, $w_x-Zq_x=w_x-Z\le0$.
\end{proof}
The engine samples that residual by inverse CDF, so a race matches it only in distribution. For a support chosen after the draft pass, retaining the top $K+1$ perturbed entries per tile tightens the maximum outside the support after up to $K$ support tokens are excluded; it is not needed for validity, since $r_i\le w_i$ makes the plain tile maximum an upper bound. A bonus sample is simpler still: a certified race over the target scores with the bonus position's field, with no partition function.

\subsection{When adaptive verification changes the sampling law}\label{app:stopping}\label{sec:stopping}
Avoiding work whose result is already determined differs from choosing which sampled proposals count. DSpark states that lossless speculation requires admission decisions that do not depend on future candidate tokens, and D-cut gives a counterexample and a repair~\cite{dspark,dcut}; the two-token illustration below is theirs in substance. Let $p=q=(1/2,1/2)$, draw $X\sim q$, discard it and draw afresh from $p$ if $X=0$, and admit it to the ordinary verifier if $X=1$. Since $p=q$ an admitted proposal is always accepted, and the output law is
\begin{equation}
\Pr(Y=0)=\tfrac14,\qquad\Pr(Y=1)=\tfrac34.\label{eq:bias}
\end{equation}
The acceptance formula was unchanged; admission selected proposals by their realized value. A mandatory progress floor on the first token does not help, since the same construction applies at the second position. Three settings avoid the issue: greedy decoding; target-only verification with point-mass proposals, where there is no proposal law to bias; and fixed-noise verification, under which every depth policy yields the same output (Theorem~\ref{thm:fixednoise}). The demand frontier below changes only the amount and order of arithmetic used to decide a fixed protocol, so it is exact pointwise.

